\documentclass[11pt,a4paper,openany]{memoir}
\newcommand{\olpath}{../upstream}
\usepackage[no-math]{fontspec}
\ifLuaTeX
\setmainfont{Nirmala UI}[Script=Tamil,Renderer=HarfBuzz,WordSpace=1.3,ItalicFont={Nirmala UI},ItalicFeatures={FakeSlant=0.15}]
\setsansfont{Nirmala UI}[Script=Tamil,Renderer=HarfBuzz,WordSpace=1.3]
\else
\setmainfont{Nirmala UI}[Script=Tamil,Renderer=OpenType,WordSpace=1.3,ItalicFont={Nirmala UI},ItalicFeatures={FakeSlant=0.15}]
\setsansfont{Nirmala UI}[Script=Tamil,Renderer=OpenType,WordSpace=1.3]
\fi
\setmonofont{Consolas}
\usepackage{xr-hyper}
\input{\olpath/sty/open-logic.sty}
\ifXeTeX
  \XeTeXlinebreaklocale "ta"
  \XeTeXlinebreakskip=0pt plus 1pt
  \XeTeXgenerateactualtext=1
\fi
\setlrmarginsandblock{25mm}{25mm}{*}
\setulmarginsandblock{22mm}{25mm}{*}
\checkandfixthelayout
\linespread{1.17}
\emergencystretch=2em
\input{../translation/tamil-config.sty}
\input{../translation/tamil-pdf-math.sty}
% Latin diacritics retain their spelling; Tamil remains in Nirmala UI.
% fontspec must be loaded with no-math so legacy mathematical accent fonts
% are not replaced by a text font lacking the mathematical bar glyph.
\newfontfamily\TALatinFont{Cambria}
\RenewDocumentCommand\L{}{{\TALatinFont\char"0141}}
\RenewDocumentCommand\l{}{{\TALatinFont\char"0142}}
\NewCommandCopy\TAOriginalAcute\'
\RenewDocumentCommand\'{m}{{\TALatinFont\TAOriginalAcute{#1}}}
\usepackage{newunicodechar}
\newunicodechar{Ł}{{\TALatinFont\char"0141}}
\newunicodechar{ł}{{\TALatinFont\char"0142}}
\newunicodechar{ś}{{\TALatinFont\char"015B}}
\AtBeginEnvironment{thebibliography}{\TALatinFont}
% cleveref's list joins and enumerate labels are generated prose; localize them
% after babel has selected the document language.
\AtBeginDocument{%
  \crefname{enumi}{படி}{படிகள்}%
  \Crefname{enumi}{படி}{படிகள்}%
  \crefname{enumii}{படி}{படிகள்}%
  \Crefname{enumii}{படி}{படிகள்}%
  \renewcommand{\crefpairconjunction}{ மற்றும் }%
  \renewcommand{\crefmiddleconjunction}{, }%
  \renewcommand{\creflastconjunction}{ மற்றும் }%
  \renewcommand{\crefpairgroupconjunction}{ மற்றும் }%
  \renewcommand{\crefmiddlegroupconjunction}{, }%
  \renewcommand{\creflastgroupconjunction}{, மற்றும் }%
  \renewcommand{\crefrangeconjunction}{--}%
}
% The upstream starred induction case and exercise marker contain hard-coded
% English even when their source arguments are translated.
\RenewDocumentCommand\indcase { s t{!} m m +m }{%
  \DeclareDocumentMacro \indfrm {#3}%
  \DeclareDocumentMacro \indfrmp {#3}%
  \DeclareDocumentMacro \indcomplex {#4}%
  \IfBooleanTF{#1}{$#3$ அணு வாய்பாடு: }{$#3 \ident #4$: }%
  \IfBooleanTF{#2}{பயிற்சி.}{#5}}

\tagtrue{novice,math,compsci,FOL,prvTrue,prvFalse,prvEx,prvAll}
\addto\captionsenglish{%
  \renewcommand{\chaptername}{அத்தியாயம்}%
  \renewcommand{\partname}{பகுதி}%
  \renewcommand{\appendixname}{இணைப்பு}%
  \renewcommand{\tablename}{அட்டவணை}%
}
\setlocalecaption{english}{photocredits}{படங்களுக்கான நன்றியுரை}
\setlocalecaption{english}{history}{வரலாற்றுக் குறிப்புகள்}
\setlocalecaption{english}{reading}{மேலும் படிக்க}
\setlocalecaption{english}{appendix}{இணைப்பு}
\setlocalecaption{english}{appendices}{இணைப்புகள்}
\setlocalecaption{english}{definition}{வரையறை}
\setlocalecaption{english}{example}{எடுத்துக்காட்டு}
\setlocalecaption{english}{lemma}{துணைத்தேற்றம்}
\setlocalecaption{english}{proposition}{கூற்று}
\setlocalecaption{english}{corollary}{தொடர்விளைவு}
\setlocalecaption{english}{problem}{பயிற்சி}
\setlocalecaption{english}{problems}{பயிற்சிகள்}
\setlocalecaption{english}{theorem}{தேற்றம்}
\setlocalecaption{english}{remark}{குறிப்பு}
\setlocalecaption{english}{axiom}{அடிகோள்}
\setlocalecaption{english}{note}{குறிப்பு}
\setlocalecaption{english}{case}{நிலை}
\setlocalecaption{english}{convention}{மரபு}
\addto\captionsenglish{\renewcommand{\proofname}{நிறுவல்}}
\addto\captionsenglish{\renewcommand{\figurename}{படம்}}
\crefname{page}{பக்கம்}{பக்கங்கள்}
\Crefname{page}{பக்கம்}{பக்கங்கள்}
\crefname{part}{பகுதி}{பகுதிகள்}
\Crefname{part}{பகுதி}{பகுதிகள்}
\crefname{chapter}{அத்தியாயம்}{அத்தியாயங்கள்}
\Crefname{chapter}{அத்தியாயம்}{அத்தியாயங்கள்}
\crefname{section}{பிரிவு}{பிரிவுகள்}
\Crefname{section}{பிரிவு}{பிரிவுகள்}
\crefname{subsection}{பிரிவு}{பிரிவுகள்}
\Crefname{subsection}{பிரிவு}{பிரிவுகள்}
\crefname{figure}{படம்}{படங்கள்}
\Crefname{figure}{படம்}{படங்கள்}
\crefname{table}{அட்டவணை}{அட்டவணைகள்}
\Crefname{table}{அட்டவணை}{அட்டவணைகள்}
\crefname{equation}{சமன்பாடு}{சமன்பாடுகள்}
\Crefname{equation}{சமன்பாடு}{சமன்பாடுகள்}
\crefname{defn}{வரையறை}{வரையறைகள்}
\Crefname{defn}{வரையறை}{வரையறைகள்}
\crefname{thm}{தேற்றம்}{தேற்றங்கள்}
\Crefname{thm}{தேற்றம்}{தேற்றங்கள்}
\crefname{ex}{எடுத்துக்காட்டு}{எடுத்துக்காட்டுகள்}
\Crefname{ex}{எடுத்துக்காட்டு}{எடுத்துக்காட்டுகள்}
\crefname{lem}{துணைத்தேற்றம்}{துணைத்தேற்றங்கள்}
\Crefname{lem}{துணைத்தேற்றம்}{துணைத்தேற்றங்கள்}
\crefname{prop}{கூற்று}{கூற்றுகள்}
\Crefname{prop}{கூற்று}{கூற்றுகள்}
\crefname{prob}{பயிற்சி}{பயிற்சிகள்}
\Crefname{prob}{பயிற்சி}{பயிற்சிகள்}
\crefname{rem}{குறிப்பு}{குறிப்புகள்}
\Crefname{rem}{குறிப்பு}{குறிப்புகள்}
\crefname{cor}{தொடர்விளைவு}{தொடர்விளைவுகள்}
\Crefname{cor}{தொடர்விளைவு}{தொடர்விளைவுகள்}
\crefname{axiom}{அடிகோள்}{அடிகோள்கள்}
\Crefname{axiom}{அடிகோள்}{அடிகோள்கள்}
\crefname{note}{குறிப்பு}{குறிப்புகள்}
\Crefname{note}{குறிப்பு}{குறிப்புகள்}
\crefname{case}{நிலை}{நிலைகள்}
\Crefname{case}{நிலை}{நிலைகள்}
\crefname{conv}{மரபு}{மரபுகள்}
\Crefname{conv}{மரபு}{மரபுகள்}
\AtBeginEnvironment{align*}{\fontsize{9.5}{12}\selectfont}
\hypersetup{pdftitle={திறந்த தருக்கவியல் — மூல மாற்று வடிவங்கள்},pdfauthor={Open Logic Project; தமிழ் மொழிபெயர்ப்புத் திட்டம்},pdfsubject={உறைந்த OpenLogic மூலத்தின் மாற்று வடிவங்கள்; இயந்திரத் தமிழ் மொழிபெயர்ப்பு}}
\addto\captionsenglish{\renewcommand{\contentsname}{பொருளடக்கம்}\renewcommand{\bibname}{மேற்கோள் நூல்}}
\tagtrue{prfSC,prfND,prfTab,prfAX}
\includeenv{editorial}
\externaldocument{tamil-complete}[tamil-complete.pdf]
\makeatletter
\newcommand\TADriverHeadingMode{\def\ol@chaptercs{section}\def\ol@partcs{section}}
\makeatother
% Alternate drivers retain their headings, editorial prose and import
% inventories. Their shared chapters are read in the principal volume.
\NewDocumentCommand\TADriverImport{s o m o}{%
  \par\noindent மூல உள்ளீடு:
  \IfBooleanT{#1}{\texttt{*}}%
  \IfNoValueF{#2}{\texttt{\detokenize{#2}/}}%
  \texttt{\detokenize{#3}.tex}%
  \IfNoValueF{#4}{\quad பிரிவு வடிவம்: \texttt{\detokenize{#4}}}%
  \par
}
\NewDocumentCommand\TACompanionUnit{m m}{%
  \clearpage
  \noindent மூல அலகு: \texttt{#1}\par
  \subfile{../translation/#2}%
}
\NewDocumentCommand\TACompanionDriver{m m m m}{%
  \clearpage
  \noindent மூல அலகு: \texttt{#1}\par
  \begingroup
  \TADriverHeadingMode
  \olchapterid{#3}{#4}%
  \RenewDocumentCommand\olimport{s o m o}{%
    \IfBooleanTF{##1}{%
      \IfNoValueTF{##2}{%
        \IfNoValueTF{##4}{\TADriverImport*{##3}}{\TADriverImport*{##3}[##4]}%
      }{%
        \IfNoValueTF{##4}{\TADriverImport*[##2]{##3}}{\TADriverImport*[##2]{##3}[##4]}%
      }%
    }{%
      \IfNoValueTF{##2}{%
        \IfNoValueTF{##4}{\TADriverImport{##3}}{\TADriverImport{##3}[##4]}%
      }{%
        \IfNoValueTF{##4}{\TADriverImport[##2]{##3}}{\TADriverImport[##2]{##3}[##4]}%
      }%
    }%
  }%
  \subfile{../translation/#2}%
  \endgroup
}
\begin{document}
\raggedright
\pagestyle{plain}
\chapter*{மூல மாற்று வடிவங்களும் கூடுதல் பகுதிகளும்}
\noindent இந்தியத் தமிழ்ப் பதிப்பு — ta-Taml-IN

\noindent மூல நூல்: Open Logic Project.
\href{https://github.com/OpenLogicProject/OpenLogic}{மூலத் திட்டம்},
\href{https://creativecommons.org/licenses/by/4.0/}{CC BY 4.0}.
மூலத் திருத்தம்: \texttt{9620cc73f9c8e0ad003c514a5d3748f29611c4c0}.

\noindent முதல் 570 அலகுகளின் இயந்திரத் தமிழாக்கமும் திருத்தங்களும்
OpenAI Codex — GPT-5.6 Sol, Ultra சிந்தனை நிலையில் செய்யப்பட்டன.
மீதமுள்ள 152 அலகுகளின் தமிழாக்கமும் பின்னைய பதிப்புத் திருத்தங்களும்
OpenAI Codex — GPT-6 Sol, Ultra சிந்தனை நிலையில் செய்யப்பட்டன.
சுயாதீனமான தாய்மொழி மதிப்புரை செய்யப்பட்டதாகக் கூறப்படவில்லை.
இந்தத் துணைநூலும் முதன்மை நூலும்
சேர்ந்து உறைந்த மூலத்தின் 722 அலகுகளையும் கணக்கிடுகின்றன.

\noindent முதன்மை வாசிப்பு நூலின் 695 மூல அலகுகளுடன்,
இங்கு 27 தனித்த மூல அலகுகள் சேர்க்கப்படுகின்றன.
அவற்றில் நான்கு மாற்று அத்தியாயம் அல்லது பகுதி அமைப்புக் கோப்புகள்.
அவற்றின் தமிழ்த் தலைப்புகள், ஆசிரியர் குறிப்புகள் மற்றும் மூல
உள்ளீட்டுப் பட்டியல்கள் இங்கே உள்ளன; ஏற்கெனவே முதன்மை நூலில்
இடம்பெறும் அத்தியாயங்கள் மீண்டும் அச்சிடப்படவில்லை.
மற்ற 23 அலகுகளின் முழுப் பாடப்பொருள் இங்கே இடம்பெறுகிறது.
அனைத்து 722 அலகுகளின் திருத்தத்தக்க மூலக் கோப்புகளும்
முழு மூலத் தொகுப்பில் உள்ளன.

\noindent \href{tamil-complete.pdf}{முதன்மை வாசிப்பு நூலைத் திறக்க}.
மூலத்தின் மாற்று வடிவங்கள் சில இடங்களில் அதே அடையாளங்களைப்
பயன்படுத்துகின்றன. இந்நூலில் உள்ள உள்ளூர் முடிவுகளின் இணைப்புகள்
இந்நூலுக்குள்ளும், பிற மேற்கோள்கள் முதன்மை நூலுக்கும் செல்கின்றன.
உறைந்த மூலத்தில் இல்லாத நிறுவல் அல்லது விளக்கம் சேர்க்கப்பட்டதாகக்
கூறப்படவில்லை.

\clearpage
\begingroup
\sloppy
\renewcommand{\cftsectionfont}{\small}
\tableofcontents*
\endgroup
\chapter{முதல்தரத் தருக்கத்தின் மாற்று வடிவங்கள்}
\TACompanionUnit{OLP-0643}{content/first-order-logic/axiomatic-deduction/provability.tex}
\TACompanionUnit{OLP-0644}{content/first-order-logic/completeness/maximally-consistent-sets.tex}
\TACompanionUnit{OLP-0645}{content/first-order-logic/syntax-and-semantics/introduction.tex}
\TACompanionDriver{OLP-0646}{content/first-order-logic/syntax-and-semantics/syntax-and-semantics.tex}{fol}{syn}

\chapter{பிற தருக்கப் பகுதிகள்}
% The definition of C precedes the theorem about representability in Q.
\TACompanionUnit{OLP-0648}{content/incompleteness/representability-in-q/c.tex}
\TACompanionUnit{OLP-0647}{content/incompleteness/representability-in-q/c-representable.tex}
\TACompanionUnit{OLP-0649}{content/intuitionistic-logic/semantics/propositions.tex}
\TACompanionUnit{OLP-0650}{content/lambda-calculus/lambda-definability/lists.tex}
\TACompanionUnit{OLP-0651}{content/lambda-calculus/syntax/conversion.tex}
\TACompanionUnit{OLP-0652}{content/many-valued-logic/sequent-calculus/proof-theoretic-notions.tex}
\TACompanionUnit{OLP-0653}{content/model-theory/basics/nonstandard-arithmetic.tex}
\TACompanionUnit{OLP-0654}{content/normal-modal-logic/syntax-and-semantics/normal-modal-logics.tex}

\chapter{நிறுவல் கோட்பாட்டின் கூடுதல் பகுதிகள்}
\clearpage
\section{வெட்டு நீக்கலின் தனித்த மூலத் துணுக்கு}
\noindent மூல அலகு: \texttt{OLP-0660}.
இது முழு அத்தியாயமாக அமைக்கப்படாத தனித்த நிறுவல் துணுக்கு.
கோப்பின் பெயர் \texttt{intuitionistic.tex} என்றாலும், அதன்
வாய்பாடுகளில் பலமுடிவு \Log{G3c} தொடரணிகளே பயன்படுத்தப்படுகின்றன.
பெயரை வைத்து நிறுவல் முறைமை மாற்றப்படவில்லை.
\begingroup
\olfileid{pt}{cut}{int}
% SEGMENT OLP-0660-S01
% Part: proof-theory
% Chapter: cut-elimination
% Auxiliary fragment: not imported into the reader

$!A \ident (!B \lif !C)$ எனில், $\pi$ பின்வருமாறு முடிகிறது:
\[
\AxiomC{}
\RightLabel{$\pi_1'$}
\Deduce$!B, \Gamma \fCenter \Delta, !C$
\RightLabel{\RightR{\lif}}
\UnaryInf$\Gamma \fCenter \Delta, !B \lif !C$
\LeftSubproofLabel{$\pi_1$}
\AxiomC{}
\RightLabel{$\pi_2'$}
\Deduce$!B \lif !C, \Pi \fCenter \Lambda, !B$
\AxiomC{}
\RightLabel{$\pi_2''$}
\Deduce$!C, \Pi \fCenter \Lambda$
\RightLabel{\LeftR{\lif}}
\BinaryInf$!B \lif !C, \Pi \fCenter \Lambda$
\RightSubproofLabel{$\pi_2$}
\RightLabel{\Cut}
\BinaryInf$\Gamma, \Pi \fCenter \Delta, \Lambda$
\DisplayProof
\]

% SEGMENT OLP-0660-S02
\Cut{} இல் முடியும் பின்வரும் !!{proof} ஐக் கருதுக:
\[
\AxiomC{}
\RightLabel{$\pi_1'$}
\Deduce$!B, \Gamma \fCenter \Delta, !C$
\RightLabel{\RightR{\lif}}
\UnaryInf$\Gamma \fCenter \Delta, !B \lif !C$
\LeftSubproofLabel{$\pi_1$}
\AxiomC{}
\RightLabel{$\pi_2'$}
\Deduce$!B \lif !C, \Pi \fCenter \Lambda, !B$
\RightLabel{\Cut}
\BinaryInf$\Gamma, \Pi \fCenter \Delta, \Lambda, !B$
\DisplayProof
\]
இதன் வெட்டு !!{height}, $\pi$ இன் வெட்டு உயரத்தைவிடக் குறைவு.
தொகுத்தறிதல் கருதுகோள் $\Gamma, \Pi \fCenter \Delta,
\Lambda, !B$ இன் !!a{proof}~$\pi_1''$ ஐத் தருகிறது.

% SEGMENT OLP-0660-S03
\Cut{} இல் முடியும் பின்வரும் !!{proof} ஐக் கருதுக:
\[
\AxiomC{}
\RightLabel{$\pi_1'$}
\Deduce$!B, \Gamma \fCenter \Delta, !C$
\AxiomC{}
\RightLabel{$\pi_2''$}
\Deduce$!C, \Pi \fCenter \Lambda$
\RightLabel{\Cut}
\BinaryInf$!B, \Gamma, \Pi \fCenter \Delta, \Lambda$
\DisplayProof
\]
இதன் வெட்டுத் தரமும் வெட்டு !!{height} உம் $\pi$ இன்
வெட்டுத் தரம், வெட்டு !!{height} ஆகிய இரண்டையும்விடக் குறைவு.
எனவே தொகுத்தறிதல் கருதுகோள் $!B, \Gamma, \Pi
\fCenter \Delta, \Lambda$ இன் \Log{G3c}-!!{proof}~$\pi_2'''$
ஐத் தருகிறது.

% SEGMENT OLP-0660-S04
இப்போது \Cut{} இல் முடியும் பின்வரும் !!{proof} க்குத்
தொகுத்தறிதல் கருதுகோளைப் பயன்படுத்துக:
\[
\AxiomC{}
\RightLabel{$\pi_1''$}
\Deduce$\Gamma, \Pi \fCenter \Delta, \Lambda, !B$
\AxiomC{}
\RightLabel{$\pi_2'''$}
\Deduce$!B, \Gamma, \Pi \fCenter \Delta, \Lambda$
\RightLabel{\Cut}
\BinaryInf$\Gamma, \Gamma, \Pi, \Pi \fCenter \Delta, \Delta, \Lambda, \Lambda$
\DisplayProof
\]
(இதன் வெட்டுத் தரம் $< \cutr{\pi}$.) ஆகவே
$\Gamma, \Gamma, \Pi, \Pi \fCenter \Delta, \Delta,
\Lambda, \Lambda$ இன் !!a{proof} கிடைக்கிறது.
\olref[seq][inv]{prop:G3c-cont-adm} ஐப் பயன்படுத்திச் சுருக்கினால்,
$\Gamma, \Pi \fCenter \Delta, \Lambda$ இன்
!!{proof}~$\pi'$ கிடைக்கிறது.

\endgroup
\TACompanionUnit{OLP-0694}{content/proof-theory/propositions-as-types/sequent-natural-deduction.tex}

\clearpage
\section{மாற்று தொடர்கணித விதி அட்டவணைகள்}
\noindent மூல அலகுகள்:
\texttt{OLP-0705}, \texttt{OLP-0708}, \texttt{OLP-0709}, \texttt{OLP-0710}.
மூலத்தில் வழங்கப்பட்ட விதிகளும் நிபந்தனைகளும் அப்படியே பாதுகாக்கப்படுகின்றன.
\begingroup
\olfileid{pt}{seq}{rai}
\subfile{../translation/content/proof-theory/sequent-calculus/rules-G1i.tex}
\clearpage
\subfile{../translation/content/proof-theory/sequent-calculus/rules-G3i.tex}
\clearpage
\subfile{../translation/content/proof-theory/sequent-calculus/rules-LK.tex}
\clearpage
\subfile{../translation/content/proof-theory/sequent-calculus/rules-mG3i.tex}
\clearpage
\endgroup

\chapter{கூற்றுத் தருக்கமும் இரண்டாம்தரத் தருக்கமும்}
\TACompanionUnit{OLP-0715}{content/propositional-logic/syntax-and-semantics/soundness.tex}
\TACompanionUnit{OLP-0714}{content/propositional-logic/syntax-and-semantics/completeness.tex}
\TACompanionUnit{OLP-0716}{content/second-order-logic/syntax-and-semantics/language-of-sol.tex}

\chapter{கணங்கள், தொடர்புகள், சார்புகளின் மாற்று வடிவங்கள்}
\TACompanionUnit{OLP-0717}{content/sets-functions-relations/functions/isomorphic-functions.tex}
\TACompanionUnit{OLP-0718}{content/sets-functions-relations/inductive-defs-proofs/introduction.tex}
\TACompanionUnit{OLP-0721}{content/sets-functions-relations/sets/proofs-about-sets.tex}
\TACompanionDriver{OLP-0719}{content/sets-functions-relations/relations/relations.tex}{sfr}{rel}
\TACompanionDriver{OLP-0720}{content/sets-functions-relations/sets-functions-relations.tex}{sfr}{}
\TACompanionDriver{OLP-0722}{content/sets-functions-relations/size-of-sets/size-of-sets.tex}{sfr}{siz}

\nocite{Frege1953}
\bibliographystyle{../upstream/bib/natbib-oup}
\bibliography{../upstream/bib/open-logic}
\end{document}

